\section{The sampled construction at successive cuts}
\label{sec:lvv_sampled_construction}

\begin{theorem}[Decomposition preserving earlier cut bounds]
    \label{thm:lvv_cut_decomposition}
    \lean{TNLean.MPS.AreaLaw.operatorSchmidtRank_cutOperator_leftConfigOfFin,
        TNLean.MPS.AreaLaw.exists_cutDecomposition_of_operatorSchmidtRank_le}
    \leanok
    Let $K$ act on a chain of $L$ sites and have operator Schmidt rank at most
    $D$ at every cut, including the endpoints. For every $i\le L$ there is a
    decomposition $K=\sum_{j=1}^{N}A_j\otimes B_j$ across the cut after $i$
    sites, with $N\le D$, such that each $A_j$ has operator Schmidt rank at
    most $D$ at every internal cut of the first $i$ sites.
\end{theorem}
\begin{proof}
    \leanok
    Choose a basis of the right matrix-block span of $K$ and extend its
    coordinate functionals to the full right operator space. Applying these
    functionals gives the left factors. A right contraction cannot increase
    the operator Schmidt rank across an earlier cut. The coordinate
    identification between an initial segment and a shorter chain preserves
    these ranks.
\end{proof}

\begin{theorem}[Simultaneous bounds for sampled factors]
    \label{thm:lvv_sampled_cut_decomposition}
    \lean{TNLean.MPS.AreaLaw.operatorSchmidtRank_cutOperator_bondComplementWord_le_all,
        TNLean.MPS.AreaLaw.exists_sampledProjector_cutDecomposition}
    \leanok
    Suppose that each letter occurs at most $\kappa$ times in each of $\ell$
    sampled words. Put $D=q^{2\kappa}$. At any cut, the sampled operator has
    a product decomposition with at most $\ell D$ terms. Each left factor
    has operator Schmidt rank at most $D$ at every earlier internal cut.
\end{theorem}
\begin{proof}
    \leanok
    \uses{thm:lvv_cut_decomposition}
    Every internal cut of a word has rank at most $q^{2\kappa}$; the
    endpoint ranks are at most one. Apply Theorem~\ref{thm:lvv_cut_decomposition}
    to each word separately. Concatenate their decompositions and absorb
    the common sampling scalar into the left factors. Scalar multiplication
    does not increase operator Schmidt rank.
\end{proof}

\begin{theorem}[Sampled error reduction]
    \label{thm:lvv_sampled_error_reduction}
    \lean{TNLean.MPS.AreaLaw.exists_sampled_errorReduction}
    \leanok
    Let $H=\sum_jH_j$ satisfy the positive nearest-neighbour contraction
    assumptions, with a unit ground vector $\Gamma$, ground energy
    $\varepsilon_0$ and positive gap $\Delta$. For $0<\delta\le1$, choose
    the sampling precision $\eta=\delta/8$, the power
    $m=\lceil((L-\varepsilon_0)/\Delta)\log(1/\eta)\rceil$, and the sample
    count $\ell$ and multiplicity bound $\kappa$ of the sampled projector.
    There is a sample with error less than $\delta/4$ from
    $|\Gamma\rangle\langle\Gamma|$. At every cut its left factors can be
    chosen before the input family, so that every $(i,s,b,1/2)$-viable
    family is sent to an $(i,\ell Ds,D\max(1,b),\delta)$-viable family,
    where $D=q^{2\kappa}$.
\end{theorem}
\begin{proof}
    \leanok
    \uses{thm:lvv_sampled_cut_decomposition, thm:lvv_mps_error_reduction}
    The power error is at most $\eta$ and the sampling error is less than
    $\eta$. Their sum is less than $\delta/4$. The simultaneous bounds of
    Theorem~\ref{thm:lvv_sampled_cut_decomposition} supply the hypotheses
    of Theorem~\ref{thm:lvv_mps_error_reduction}.
\end{proof}

\begin{theorem}[Uniform polynomial sampling bounds]
    \label{thm:lvv_sampled_polynomial_bounds}
    \lean{TNLean.MPS.AreaLaw.exists_sampledProjectorWordCount_polynomial,
        TNLean.MPS.AreaLaw.exists_sampledProjectorLetterBound_logarithmic,
        TNLean.MPS.AreaLaw.exists_sampledProjector_uniform_polynomial}
    \leanok
    Fix $q\ge1$, $\Delta>0$, $0<c\le1$ and a real bound $B$.
    There are constants $C\ge1$ and $p\in\mathbb N$, chosen before $L$
    and $\varepsilon_0$, such that whenever $L\ge1$,
    $0\le\varepsilon_0\le L-1$ and $\varepsilon_0/\Delta\le B$, sampling
    at precision $c/L$ gives $\ell D\le CL^p$ and $D\le CL^p$.
    The count $\ell$ has a polynomial bound and $\kappa$ a logarithmic bound.
\end{theorem}
\begin{proof}
    \leanok
    Set $r=\lceil B\rceil$. The normalization obeys
    $C_0^m\le L(L/c)^r$. Substitution in the sample-count formula gives
    $\ell\le C_1L^{2r+5}$. Also
    $m/L\le\Delta^{-1}(\log(1/c)+\log L)+1$.
    Substitution in $\kappa=\lceil2m/L+2\log(4\ell L^4)\rceil$ gives
    $\kappa\le a\log L+b$. Thus
    $q^{2\kappa}\le\exp(2b\log q)L^{\lceil2a\log q\rceil}$.
    Multiplying these bounds and rounding the coefficient proves the claim.
\end{proof}

\begin{definition}[Successive viability precision]
    \label{def:lvv_successive_precision}
    \lean{TNLean.MPS.AreaLaw.lvvViabilityPrecision}
    \leanok
    For a positive gap $\Delta$ and target accuracy $\eta>0$, set
    $\alpha=\min\{1,\eta,(\Delta/169)^2\}$.
\end{definition}

\begin{theorem}[A complete step with fixed polynomial bounds]
    \label{thm:lvv_sampled_successor}
    \lean{TNLean.MPS.AreaLaw.exists_sampled_viable_successor_polynomial}
    \leanok
    \uses{def:lvv_successive_precision}
    Fix $q\ge2$, $\Delta>0$, $\eta>0$ and an upper bound $B$ for
    $\varepsilon_0/\Delta$. There are natural constants $C,D\ge1$ and
    $p,r$ such that, for every positive chain length and every cut
    $1\le i\le L$, any finite family of left matrix product states
    that is $(i-1,\alpha/L)$-viable gives an
    $(i,CL^p,DL^r,\alpha/L)$-viable family by extension, size trimming,
    independent bond trimming and sampled error reduction.
    The bounds are independent of the size and bond bound of the input family.
\end{theorem}
\begin{proof}
    \leanok
    \uses{thm:lvv_sampled_error_reduction, thm:lvv_sampled_polynomial_bounds,
        thm:lvv_size_bond_trimmed_family, lem:lvv_bond_trimming_radius_polynomial}
    Extension preserves viability. Since $\alpha\le(\Delta/169)^2$,
    size and bond trimming produce a half-viable family of polynomial
    cardinality and individual bond bound. Sampled error reduction restores
    error $\alpha/L$. Multiply its polynomial factor counts and ranks by
    the corresponding trimming bounds. All coefficients and degrees are
    chosen before $L$.
\end{proof}

\begin{theorem}[Successive viable families]
    \label{thm:lvv_sampled_successive_families}
    \lean{TNLean.MPS.AreaLaw.exists_sampled_successive_viable_sets_polynomial}
    \leanok
    Under the hypotheses of Theorem~\ref{thm:lvv_sampled_successor}, there
    are families $S_i$ at all cuts $0\le i\le L$, beginning with the scalar
    singleton, such that each is $(i,CL^p,DL^r,\alpha/L)$-viable.
\end{theorem}
\begin{proof}
    \leanok
    \uses{thm:lvv_sampled_successor, lem:lvv_initial_scalar_family}
    The scalar singleton is exactly viable. Apply
    Theorem~\ref{thm:lvv_sampled_successor} inductively at the successive
    cuts. The same size and bond bounds apply at each step.
\end{proof}

\begin{theorem}[A matrix product state in the final span]
    \label{thm:lvv_final_span_mps}
    \lean{TNLean.MPS.AreaLaw.isOpenMPS_leftChain_of_mem_span,
        TNLean.MPS.AreaLaw.IsViable.exists_openMPS_of_forall_isLeftOpenMPS,
        TNLean.MPS.AreaLaw.IsViableMPS.exists_openMPS}
    \leanok
    A final $(L,s,b,\delta)$-viable family contains in its span a normalized
    matrix product state of bond at most $s\max(1,b)$ and overlap at least
    $1-\delta$ with the ground vector.
\end{theorem}
\begin{proof}
    \leanok
    \uses{lem:lvv_final_span_vector}
    At the final cut the viable witness belongs to the span of the family.
    Matrix rank is subadditive under finite sums, so its internal cut ranks
    are at most the sum of the individual bond bounds.
\end{proof}

\begin{theorem}[Polynomial final bond at prescribed accuracy]
    \label{thm:lvv_sampled_ground_state}
    \lean{TNLean.MPS.AreaLaw.exists_sampled_groundState_openMPS_polynomial}
    \leanok
    For fixed $q\ge2$, $\Delta>0$, $\eta>0$ and $B$, there are natural
    constants $C\ge1,p$ such that every positive nearest-neighbour
    contraction chain of length $L\ge1$, with a unique gapped ground vector
    $\Gamma$ and $\varepsilon_0/\Delta\le B$, admits a normalized matrix
    product state $\psi$ of bond at most $CL^p$ with
    $|\langle\psi,\Gamma\rangle|\ge1-\eta$.
    The construction is the successive sampled construction.
\end{theorem}
\begin{proof}
    \leanok
    \uses{thm:lvv_sampled_successive_families, thm:lvv_final_span_mps}
    Apply Theorem~\ref{thm:lvv_final_span_mps} to $S_L$.
    Multiply the size and bond polynomials. The overlap error is at most
    $\alpha/L\le\eta$.
\end{proof}
